\chapter{2006 Nobel Prize in Economics}
\textbf{Laureate: }Edmund Phelps (USA)

\section*{Reason for Award}
For his research into intertemporal tradeoffs in macroeconomic policy.

\section{What They Did}
Edmund Phelps made two major contributions to macroeconomic theory that reshaped our understanding of inflation, unemployment, and the limits of stabilization policy. His work addressed fundamental questions about the relationship between short-run policy actions and long-run economic outcomes.

Phelps developed the natural rate of unemployment hypothesis, independently and contemporaneously with Milton Friedman, demonstrating that the Phillips curve trade-off between inflation and unemployment exists only in the short run. In his seminal 1967 paper, ``The Role of Monetary Policy,'' Phelps argued that any attempt to push unemployment below its natural rate through expansionary monetary policy would only produce temporary gains, as workers and firms adjust their inflation expectations upward. Once expectations adjust, unemployment returns to its natural rate, but with higher inflation.

The natural rate is determined by structural factors in the economy---labor market institutions, education and training systems, demographic patterns, technological change, and the degree of competition in product and labor markets. It is not a fixed constant but can shift over time as these underlying factors change. Phelps' insight was that monetary policy can affect unemployment only in the short run, when expectations are slow to adjust, but cannot influence it in the long run. This insight placed fundamental limits on the ability of monetary policy to stimulate employment.

Phelps developed the expectations-augmented Phillips curve, which incorporates inflation expectations into the relationship between inflation and unemployment. The equation takes the form: inflation equals expected inflation minus a coefficient times unemployment minus the natural rate, plus supply shocks. This formulation showed that the trade-off disappears once expected inflation adjusts to actual inflation, leaving only the effect of unanticipated inflation on unemployment. The model provided a rigorous foundation for understanding why attempts to maintain unemployment below the natural rate through persistent inflationary policy ultimately fail.

Phelps also developed the theory of intertemporal macroeconomic policy, examining how current decisions affect future welfare. His work on consumption-saving decisions across lifetimes, building on the lifecycle hypothesis, showed how individuals smooth consumption over time based on expected lifetime resources. His analysis of investment choices and their effects on future productive capacity illuminated the trade-offs between current consumption and future growth. Phelps examined the long-run implications of monetary and fiscal policy, showing that policies designed to provide temporary stimulus can have persistent effects on the structure of the economy if they alter expectations, investment incentives, or institutional arrangements.

\section{Impact on Economic Thinking}
Phelps' natural rate hypothesis fundamentally transformed macroeconomic understanding of the inflation-unemployment relationship. It demonstrated that monetary policy cannot permanently reduce unemployment below the natural rate; stimulating demand merely accelerates inflation without lasting employment gains. This insight ended the belief, held by many Keynesian economists in the 1960s, that policymakers could permanently trade off higher inflation for lower unemployment.

The distinction between cyclical unemployment (caused by insufficient demand) and structural unemployment (caused by mismatches between worker skills and job requirements, or by institutional rigidities) became crucial for policy formulation. Policies aimed at reducing cyclical unemployment through demand stimulus are appropriate during recessions, but structural unemployment requires supply-side interventions such as education, training, labor market flexibility, and removal of regulatory barriers. This distinction has guided policy debates about the appropriate mix of monetary and structural policies.

Central banks adopted the natural rate concept as a cornerstone of macroeconomic analysis. The understanding that there is no long-run trade-off between inflation and unemployment led to the prioritization of price stability as the primary goal of monetary policy. The expectations-augmented Phillips curve became standard in economics textbooks, clarifying the limits of accommodative monetary policy and emphasizing the importance of anchoring inflation expectations.

Phelps' work on intertemporal optimization influenced the development of dynamic macroeconomic models and the lifecycle theory of consumption. His insights about the time consistency of policy complemented Kydland and Prescott's work, reinforcing the importance of credible commitment in macroeconomic management.

\section{Modern Perspectives Today}
The natural rate hypothesis continues to inform central bank forward guidance and communication strategies. Modern monetary policy frameworks anchor expectations about future policy paths, preventing destabilizing speculation and reducing the cost of disinflation. The concept of the natural rate of unemployment remains central to central bank analysis, with estimates of the non-accelerating inflation rate of unemployment (NAIRU) guiding monetary easing and tightening decisions.

The natural rate is not static. Contemporary research explores how structural factors---demographic changes, technological disruption, globalization, the rise of the gig economy, occupational licensing reforms, and changes in labor market institutions---shift the natural rate over time. Understanding these dynamics is essential for calibrating monetary policy and designing effective structural reforms. The debate about whether the natural rate has declined in recent decades illustrates the ongoing relevance of Phelps' framework.

Fiscal policy analysis examines the sustainability implications of borrowing and spending, recognizing that current deficits constrain future policy options. The concept of intertemporal budget constraints, rooted in Phelps' work, reminds policymakers that today's choices affect tomorrow's possibilities. Public debt sustainability analysis builds on the intertemporal framework that Phelps helped develop.

Phelps' framework remains essential for understanding long-run economic performance and the limitations of short-term stabilization policy. His work reminds us that sustainable prosperity requires investment in human capital, structural reform, and policies that enhance productive capacity, not just demand management. The natural rate concept continues to shape debates about the appropriate role of monetary policy, the design of fiscal rules, and the importance of structural reform in addressing persistent unemployment.
